% !TEX root = ../main.tex
%-------------------------------------------------------------------------------
%	ACKNOWLEDGEMENTS
%-------------------------------------------------------------------------------

\begin{acknowledgements}
	\addchaptertocentry{\acknowledgementname}
	I would like to express my deepest gratitude to those who supported me
	throughout this journey.

	First and foremost, I sincerely thank my supervisor, \supname, for the
	continuous guidance, support, and encouragement
	throughout this work.
	The opportunity to carry out this research under his supervision and to
	gain such a valuable and unforgettable research experience made this year
	truly special for me.

	I would also like to express my sincere gratitude to my examiner,
	\examname, for making this master's thesis possible in the first place.
	By establishing contact with my supervisor, placing her trust in me, and
	supporting the organisational aspects throughout the process, she made both
	this work and my year abroad possible.
	I am deeply grateful for her support.

	Special thanks go to \'Alvaro Tejero Moyano, Rhea Alexander, Gustavo
	Mereles Menesse, Giulio Camillo da Silva, and the other members of the
	research group for their kindness, helpfulness, and the pleasant working
	atmosphere.
	Their support made this experience both productive and enjoyable.

	Finally, I would like to express my heartfelt gratitude to my girlfriend,
	Anna Schupeta, for her constant love, patience, and unwavering support throughout
	this entire process.
	She supported me tremendously through all the highs and lows, both
	emotionally and academically, and I am deeply grateful for that.
	I am also very thankful to my mother, whose outside perspective and
	thoughtful advice often provided valuable support.

\end{acknowledgements}

%-------------------------------------------------------------------------------
%	LIST OF CONTENTS/FIGURES/TABLES PAGES
%-------------------------------------------------------------------------------

\tableofcontents % Prints the main table of contents

%\listoffigures % Prints the list of figures
%\listoftables % Prints the list of tables

\clearpage
%-------------------------------------------------------------------------------
%	NOTATION AND CONVENTIONS
%-------------------------------------------------------------------------------
\subsection*{Notation and conventions}

Throughout this thesis, natural units are used with $\hbar = 1$ and $c = 1$.
Frequencies, energies, and inverse times are expressed in the same units.

The following notation conventions apply:

\begin{itemize}

\item \textbf{Vectors.}
Vectors in real space are denoted by bold symbols, e.g.\ $\boldsymbol{E}$ (electric field), $\boldsymbol{P}$ (polarisation), $\boldsymbol{k}$ (wave vector), $\boldsymbol{d}$ (transition dipole moment), and $\boldsymbol{e}$ (polarisation unit vector). 
The magnitude of a vector $\boldsymbol{d}$ is written as $d = |\boldsymbol{d}|$ when required.

\item \textbf{Operators and states.}
Operators acting on Hilbert spaces carry a hat, e.g.\ $\hat{H}$ for the Hamiltonian and $\hat{\boldsymbol{\mu}}$ for the dipole operator. 
Density operators are written as $\hat{\rho}$. 
Matrix elements in a chosen basis are denoted by plain symbols, e.g.\ $\rho_{ij} = \langle i | \hat{\rho} | j \rangle$.

\item \textbf{Quantum mechanical pictures.}
Different dynamical representations (Schr\"odinger, interaction, Heisenberg) are indicated by subscripts. 
For example, $\hat{O}_I(t)$ denotes an operator in the interaction picture. 
To avoid ambiguity with this convention, the interaction Hamiltonian is written as $\hat{H}_{\mathrm{int}}$.

\item \textbf{Superoperators.}
Linear maps acting on operators (e.g.\ propagators or Liouvillians) are written in calligraphic font, such as $\mathcal{U}$, $\mathcal{L}$. 
Their action on density operators is written explicitly as $\mathcal{L}[\hat{\rho}]$.

\item \textbf{Field projections.}
For a fixed laboratory polarisation direction $\boldsymbol{e}$, vector fields are often projected onto this axis. 
The scalar field envelope is then defined as
\[
E(t) = \boldsymbol{e} \cdot \boldsymbol{E}(t),
\]
and correspondingly $\hat{\mu} = \hat{\boldsymbol{\mu}} \cdot \boldsymbol{e}$. 
Whenever such projections are used, the underlying vector character is understood but suppressed in the notation.

\item \textbf{Complex quantities.}
The symbols $\Re$ and $\Im$ denote the real and imaginary parts of a complex quantity, respectively.

\end{itemize}

%-------------------------------------------------------------------------------
%	ABBREVIATIONS
%-------------------------------------------------------------------------------

\iffalse
	\begin{abbreviations}{ll} % Include a list of abbreviations (a table of two columns)

		\textbf{2DES} & \textbf{2D}imensional \textbf{E}lectronic \textbf{S}pectroscopy\\
		\textbf{NIR} & \textbf{N}ear \textbf{I}nfrared\\
		\textbf{IR} & \textbf{I}nfrared\\
		\textbf{UV} & \textbf{U}ltraviolet\\
		\textbf{NP} & \textbf{N}on\text{P}erturbative\\
		\textbf{FWM} & \textbf{F}our \textbf{W}ave \textbf{M}ixing\\
		\textbf{DM} & \textbf{D}ensity \textbf{M}atrix\\

	\end{abbreviations}
\fi

%-------------------------------------------------------------------------------
%	PHYSICAL CONSTANTS/OTHER DEFINITIONS
%-------------------------------------------------------------------------------

\iffalse
	\begin{constants}{lr@{${}={}$}l} % The list of physical constants is a three column table

		Constant Name & $Symbol$ & $Constant Value$ with units\\
		Speed of Light & $c$ & $1$\\
		Planck's Constant & $\hbar$ & $1$\\

	\end{constants}
\fi

%-------------------------------------------------------------------------------
%	SYMBOLS
%-------------------------------------------------------------------------------

\iffalse
	\begin{symbols}{lll} % Include a list of Symbols (a three column table)

		$a$ & distance & natural-unit length scale \\
		$P$ & polarisation amplitude & natural units \\
		%Symbol & Name & Unit \\

		\addlinespace % Gap to separate the Roman symbols from the Greek

		$\omega$ & angular frequency & inverse time (natural units) \\

	\end{symbols}
\fi


%-------------------------------------------------------------------------------
%	DEDICATION
%-------------------------------------------------------------------------------
\iffalse
	\dedicatory{For/Dedicated to/To my\ldots}
\fi
